\scanpage{38}
Soit $G$ un schéma en groupes simple sur $S=\operatorname{Spec}A$, $A$ Artinien loc.
\[
\mathfrak g=L_S(G)
\]

(1) Si $H^0(\mathfrak g_0,\mathfrak g_0)=0$, $\Longleftrightarrow$ $G\to\underline{\operatorname{Aut}}_S(\mathfrak g)$ est un \unsure{m. infi.}
[donc si $G_0$ n'a pas de sous-groupe \deleted{discret} étale \deleted{\unsure{fini invariant}} non trivial, alors c'est une \unsure{immersion} fermée]

(2) \deleted{Si} \inserted{(1) et} \deleted{de plus} $H^1(\mathfrak g_0,\mathfrak g_0)=0$ $\Longleftrightarrow$ \deleted{\unsure{donc}} $G\to\underline{\operatorname{Aut}}_S(\mathfrak g)$ est étale,
[alors sous la condition dite plus haut, c'est une immers. ouverte, donc $G\simeq\underline{\operatorname{Aut}}_S^0(\mathfrak g)$ si $G$ à fibres connexes]

\margin{\deleted{(1) et (2)} \unsure{(2)} $\Rightarrow\underline{\operatorname{Aut}}_S(\mathfrak g)\simeq\underline{\operatorname{Aut}}_S(G)$ si $G$ connexe\ldots\ Utilisant \unsure{Tôhoku} cela suffirait à construire les \unsure{Tchevalley} \deleted{\unsure{sur}} sur $\mathbf Z$\ldots}

Supposons qu'on ne dispose pas de $G$, mais seulement de $\mathfrak g$, considérons $H=\underline{\operatorname{Aut}}_S(\mathfrak g)$.
\deleted{Alors} C'est un schéma en groupes sur $S$, et $\operatorname{Lie}_k(H_0)=\mathfrak h_0$ est isomorphe à \deleted{$\unsure{\mathfrak g_0/(\text{centre de }\mathfrak g_0)}$ \unsure{par}} \inserted{l'algèbre de Lie des} \deleted{\illegible{2}} dérivations de $\mathfrak g_0$, d'où $\mathfrak g_0\to\mathfrak h_0$, qui est un \unsure{isom} si (1) et (2) sont satisfaites.

S'il en est ainsi, on ne peut affirmer encore que $H$ est simple sur $S$, \unsure{avec les hyp 1 2 seules citées}, \deleted{\unsure{Mais donc}} on le peut si on suppose
\[
\text{(3)}\qquad H^2(\mathfrak g_0,\mathfrak g_0)=0.
\]
qui signifie sauf erreur qu'il n'y a pas de variations de structure \inserted{infinitésimales} \unsure{non triviales} de $\mathfrak g_0$, et \uline{implique} que \deleted{\unsure{$H$ est donc}} $\underline{\operatorname{Aut}}_S(\mathfrak g)$ est simple sur $S$.
\margin{37}
